% !TEX root = ../main.tex

% 结论
\begin{conclusions}
蛋白质功能预测旨在根据已有的生物学信息推断蛋白质可能具备的功能属性，是连接蛋白质序列数据与生命科学应用的重要桥梁。面对蛋白质数据规模持续增长而人工实验注释相对滞后的现实矛盾，如何构建具有良好泛化能力、能够适应复杂标签体系并有效利用多源信息的预测模型，已成为该领域的重要研究方向。尤其是在基因本体标签不断扩展、类别分布高度不均衡的背景下，传统将功能预测视为固定标签集分类问题的方法，逐渐暴露出标签语义建模不足、对未见功能标签适应性有限等问题。围绕上述挑战，本文从开放标签预测与参数高效适配两个层面出发，对基于多模态预训练的蛋白质功能预测方法进行了系统研究。

第一，提出了多模态零样本蛋白质功能注释方法MZSGO。该方法突破了传统封闭分类范式的限制，将蛋白质功能预测重新建模为蛋白质与GO术语之间的语义匹配任务。具体而言，该方法分别利用蛋白质语言模型和文本嵌入模型对序列进化信息、结构域描述语义以及GO功能定义进行编码，并结合自适应门控融合机制与非对称特征丢弃策略，实现多源异构信息在统一语义空间中的有效融合。实验结果表明，MZSGO不仅在分子功能、细胞组分和生物过程三个GO子本体的有监督预测任务中表现出较强竞争力，而且在零样本预测场景下展现出更为显著的优势。其中，在细胞组分子本体上的调和均值达到0.6669，较对比方法有显著提升，验证了基于语义匹配的开放标签预测范式在蛋白质功能注释任务中的有效性。

第二，提出了基于域感知适配器预训练的蛋白质功能预测方法MZSGO-DA。该方法在蛋白质语言模型的Transformer层中引入轻量级适配器模块，通过“序列—结构域”和“序列—功能”双通道对比学习，将结构域相关的功能语义以参数化形式注入适配器之中。在下游任务阶段，该方法采用冻结适配器与可训练适配器并行的双适配器设计，在保留预训练知识的同时增强了模型对具体功能预测任务的适应能力。实验结果表明，MZSGO-DA在推理阶段仅需输入蛋白质序列，便可在分子功能子本体上取得优于MZSGO的结果，在其余子本体上也取得了可比甚至更优的综合表现；同时，该方法仅需训练约7.3\%的参数即可实现有效适配，体现出较好的参数效率和实际部署潜力。

综合全文研究可以看出，充分挖掘蛋白质序列表示、结构域语义信息与基因本体标签语义之间的内在联系，是提升蛋白质功能预测性能的重要途径。本文从显式多模态语义融合到隐式语义知识注入两个层面展开研究，说明了语义信息不仅能够改善模型对已知功能标签的判别能力，也能够增强其在开放标签场景中的泛化能力。本文研究表明，蛋白质功能预测不应仅停留于传统封闭分类框架，而应进一步向面向语义关系建模和可扩展预测的方向发展。

展望未来，本论文的研究仍可从以下几个方向进一步拓展：一是将蛋白质三维结构信息进一步纳入统一的多模态建模框架，以丰富功能表征的结构层面信息；二是探索GO本体图结构约束与标签语义嵌入的联合建模，以增强标签预测结果在层级关系上的一致性；三是推动序列编码器与文本编码器的联合端到端优化，以进一步释放跨模态协同学习的潜力。上述方向有望进一步提升蛋白质功能预测的精度、鲁棒性与泛化能力，推动计算方法在大规模蛋白质功能解析中的深入应用。

\end{conclusions}